\clearpage
\section{บทสรุป ข้อค้นพบข้ามภารกิจ และแนวทางการพัฒนาต่อ (Conclusion \& Future Work)}
\label{sec:conclusion}

\subsection{บทเรียนสำคัญจากการสร้างและประเมินโมเดล (Lessons Learned Across Tasks 1, 2, 3 \& 5)}
\label{subsec:lessons_learned}

จากการพัฒนาแบบจำลองและตรวจสอบกระบวนการสร้างข้อมูลใน Task 1, Task 2, Task 3 และ Task 5 (ส่วน Task 4 ยังคงสงวนไว้ในหัวข้อย่อย~\ref{subsec:task4} ตามข้อตกลง) ทีม ML003 ได้สังเคราะห์บทเรียนเชิงวิศวกรรมข้ามภารกิจที่สำคัญที่สุด 5 ประการ ซึ่งสามารถนำไปประยุกต์ใช้กับโจทย์ Machine Learning เชิงอุตสาหกรรมได้ทั่วไป ดังนี้

\begin{enumerate}
  \item \textbf{Metric-First Decision Rules ให้ผลตอบแทนสูงกว่าการปรับจูนสถาปัตยกรรมที่ซับซ้อน:}
  การแยกชั้นความน่าจะเป็น (Ranking/Calibration) ออกจากกฎการตัดสินใจขั้นสุดท้าย (Decision Rule) ที่อนุพันธ์ตรงตามตัวชี้วัดเป้าหมาย สร้างผลลัพธ์เชิงประจักษ์ที่เหนือกว่าการเปลี่ยนโมเดลหลายเท่า:
  \begin{itemize}
    \item \emph{Task 1 (Binary-$F_2$):} การพิสูจน์ว่าเกณฑ์ประเมินคือ Binary-$F_2$ (ปรับ $FN$ หนักกว่า $FP$ $4$ เท่า) แทน Macro-$F_2$ ทำให้จุดทำงานเปลี่ยนจากอัตราบวก $\sim 41\%$ เป็น $62.40\%$ บน OOF ($\text{thr}=0.3769$; $59.33\%$ บนชุดทดสอบ) ส่งผลให้คะแนน Public LB เพิ่มขึ้นถึง $\mathbf{+0.01276}$ ($0.66614 \to 0.67890$ ใน \feat{probe\_v10e})
    \item \emph{Task 2 (Macro-$F_1$):} การแทนที่ $\arg\max_k p_k(x)$ ด้วยกฎ Fixed-Point Affine $F_1$ ($a_k = \frac{2}{\pi_k+q_k}$, $b_k = \frac{-F_k}{\pi_k+q_k}$) ช่วยชดเชยน้ำหนักคลาสส่วนน้อยบริเวณรอยต่อ ปรับเปลี่ยน 2,056 แถวบน OOF ($1.98\%$) เพิ่ม OOF Macro-$F_1$ จาก $0.75798$ เป็น $\mathbf{0.75826}$ และเพิ่ม LB $+0.00034$ ใน \feat{sub\_v18a}
    \item \emph{Task 3 (NMAE / L1):} การพิสูจน์ทฤษฎี Hurdle Mixture Median ($q = 1 - \frac{1}{2p}$ เมื่อ $p > 0.5$ และทำนาย $0$ เมื่อ $p \le 0.5$) จำกัดให้แบบจำลองควอนไทล์เรียนรู้เฉพาะช่วง $\alpha \in [0.02, 0.50]$ หลีกเลี่ยงอคติเชิงบวกของการใช้ค่าคาดหมายบนข้อมูลที่เบ้ขวาอย่างรุนแรง
    \item \emph{Task 5 ($F_3$):} การหาจุดคุ้มทุนทางคณิตศาสตร์ $p^* = \frac{F_3}{1+\beta^2} = \frac{F_3}{10} \approx 0.10$ เป็นตัวกำหนดจุดตัดที่แม่นยำสำหรับ Stage D NaN-Aware Naive Bayes ($\text{posterior} > 0.10$, กู้คืน $+52$ แถว) ขณะเดียวกันก็อธิบายเหตุผลที่การเติมแถวหางที่มีความมั่นใจต่ำกว่า 10\% (\feat{sub\_v87e} $+11$ แถว หรือ \feat{loose} $3{,}602$ แถว) ทำให้คะแนน LB ลดลงทุกครั้ง
  \end{itemize}

  \item \textbf{การแกะรอยกระบวนการสร้างข้อมูล (DGP) เป็นตัวกำหนดตระกูลโมเดลที่เหมาะสมที่สุด:}
  ไม่มีอัลกอริทึมเดียวที่ชนะทุกโจทย์ ตารางที่~\ref{tab:cross_task_dgp_summary} สรุปความเชื่อมโยงระหว่างโครงสร้าง DGP กับการเลือกสถาปัตยกรรมที่ถูกต้องในทั้ง 4 ภารกิจ:
  \begin{itemize}
    \item ใน \emph{Task 1} ขอบเขตกลุ่มย่อยแบบกล่องสี่เหลี่ยมตั้งฉาก (Uniform Boxes) บวกกับ MCAR NaNs ทำให้ GBDT พร้อม Bayesian Zone Posteriors (\feat{fe\_v2}) และ MCAR 4x Masking โดดเด่นกว่าโครงข่าย MLP ($0.88548$ AUC) และ Mixture NB ($0.86957$ AUC)
    \item ใน \emph{Task 2} ค่าสหสัมพันธ์ภายในคลาสที่ลดลง $9.6\times$ เหลือ $|\bar{r}| = 0.0036$ ($\max |r| = 0.0130$) ร่วมกับการสลับป้ายกำกับสมมาตร $3\%$ พิสูจน์ว่า Histogram EM-NB เป็นโมเดลที่ปราศจากอคติเชิงโครงสร้าง ($\text{Gap} = +0.0016$ เทียบกับ LightGBM ที่โอเวอร์ฟิตถึง $+0.0416$)
    \item ใน \emph{Task 3} กลไกตัวคูณความชันหนึ่งหน่วยจาก \feat{prev\_month\_spending} พร้อม Step Factors ไม่กี่ตัวทำให้ CatBoost \feat{SymmetricTree} ตื้น (\feat{depth=4}) บน Residual $r = \ln(y) - \text{offset}$ ทำคะแนน $0.79505$ (เดี่ยว) และ $0.79537$ ($10\times 10$ Bag) ชนะต้นไม้ลึก ($\ge 5$), Fitted Linear/Ridge Offsets (ที่ถูก Whale Spikes บิดเบือน) และ MLP ($0.78430$)
    \item ใน \emph{Task 5} พื้นหลังปกติที่มีความเป็นอิสระ ($\max |r| = 0.0263$) ผสมกับกล่อง 1--3 มิติขอบคมทำให้วงจรสกัดส่วนเหลือและ Support Box (Est.\ $F_3 \approx 0.99681$) เหนือกว่า Isolation Forest (FP นอกชุด $42.2\%$) และ Autoencoder TMAE (FP $82.8\%$) ซึ่งถูกความแปรปรวนของฟีเจอร์ปกติกว่า 20 ตัวกลบสัญญาณ
  \end{itemize}

  \item \textbf{Nested Validation และวินัยต่อผล Negative A/B ป้องกันภาพลวงตาของการจูน:}
  ตลอดโครงการ การยึดถือ Nested CV และ Paired Fold Comparisons ช่วยป้องกันการตัดสินใจผิดพลาดซ้ำซาก ไม่ว่าจะเป็นการจับได้ว่า Optuna บน Macro-$F_1$ ของ Task 2 โอเวอร์ฟิตจาก $0.7557$ (in-CV) ลงมาเหลือ $0.7544$ (Nested, ต่ำกว่า Untuned $0.7550$), การตรวจพบความลำเอียงของ In-sample blending ใน Task 3 ($0.79195$ vs Honest $0.79191$), และการกล้าปฏิเสธสมมติฐานที่ดูน่าเชื่อถือแต่ไม่เพิ่มคะแนนจริงบน Pipeline ปลายทาง (เช่น กฎ \feat{is\_premium\_member = (vip\_status >= 3)} ใน Task 3 ที่ลด log-residual บนแถวหายได้ $-4\%$ แต่ลด NMAE ของระบบจริงลง $-0.00007$; การเพิ่ม Optuna รอบสองใน Task 1; และการเพิ่ม 11 แถวขอบใน Task 5 ที่ลด LB ลง $-0.00010$)

  \item \textbf{การวัดเพดานเชิงปริมาณของข้อมูล (Data Ceiling) ช่วยให้รู้จุดหยุดที่เหมาะสม:}
  การฟิต Power-Law Learning Curve และวิเคราะห์ Error Budget ในทุกภารกิจช่วยให้ทีมทราบขอบเขตที่ลดทอนไม่ได้ (Irreducible Noise) ของข้อมูล ได้แก่ Task 1 มีเส้นกำกับ $\text{AUC} = 0.8939$ / $\text{Macro-}F_2 = 0.7833$ (โมเดลสุดท้ายทำได้ $0.89499$ / $0.78279$); Task 2 มีสัญญาณรบกวนป้ายกำกับ $3\%$ และการเติบโต $< +0.0018$ ต่อการเพิ่มข้อมูลสองเท่า (แตะเพดาน $0.758\text{--}0.760$ ที่ $0.75826$); Task 3 มีกว่า $97\%$ ของความคลาดเคลื่อนมาจากสัญญาณรบกวนสุ่มของตัวกำเนิดข้อมูล ($49.0\%$ Trapezoid noise, $32.9\%$ coin-flip ของ $y=0$, $15.4\%$ Whale Spikes 266 แถว) ทำให้เพดานอยู่ที่ $\sim 0.796\text{--}0.797$ (เทียบ $0.79537$); และ Task 5 มีส่วนเหลือสะอาด 23,785 แถวที่ผ่าน Rank-Grid Scan ด้วยค่า $z_{\max} = 6.27$ อยู่ในช่วงสุ่ม Null ($5.48\text{--}6.62$) การรู้เพดานเหล่านี้ช่วยให้ทีมหยุดวิศวกรรมฟีเจอร์ที่เกินจำเป็นแล้วหันมาลดความแปรปรวนด้วย Seed Bagging และทำให้โค้ดปลอด Leakage

  \item \textbf{ความซื่อสัตย์แบบ Test-Free Inference และ Reproducibility เป็นหัวใจของระบบจริง:}
  การปรับ Task 1 สู่ OOF ECDF + Fixed Threshold, การใช้ Train-Only Imputer ใน Task 3, และการรักษา Pipeline ที่ไม่พึ่งพารหัส \feat{player\_id} ใน Task 2 และ Task 5 ทำให้แบบจำลองทุกตัวพร้อมรองรับการทำนายรายแถวแบบ Streaming ในสภาพแวดล้อมจริง (Production) โดยไม่ขึ้นกับขนาดหรือองค์ประกอบของชุดทดสอบที่ส่งเข้ามาพร้อมกัน
\end{enumerate}

\begin{table}[htbp]
  \centering
  \caption{ตารางสังเคราะห์ข้ามภารกิจ: โครงสร้าง DGP กับการเลือกตระกูลโมเดลและบทเรียนสำคัญ}
  \label{tab:cross_task_dgp_summary}
  \footnotesize
  \setlength{\tabcolsep}{4pt}
  \begin{tabularx}{\textwidth}{@{} >{\raggedright\arraybackslash}p{1.6cm} >{\raggedright\arraybackslash}p{3.9cm} >{\raggedright\arraybackslash}p{3.8cm} X @{}}
    \toprule
    \textbf{ภารกิจ} & \textbf{กลไก DGP ที่ค้นพบ} & \textbf{ตระกูลโมเดลที่เหมาะสมที่สุด} & \textbf{เหตุผลที่ทางเลือกอื่นด้อยกว่าและเพดานข้อมูล} \\
    \midrule
    \textbf{Task 1} & 5 Uniform Boxes (34.2\% โซนบริสุทธิ์) + Gaussian Skill Mean-Shift + เอกลักษณ์ \feat{lps} + MCAR NaNs & GBDT (\feat{fe\_v2} + MCAR 4x) Blend + Train-Only ECDF Cut ($\text{thr}=0.3769$) & MLP ($0.88548$) และ NB ($0.86957$) อ่อนต่อขอบเขตสี่เหลี่ยมและสหสัมพันธ์ทักษะ $\rho \approx 0.3$; แตะเพดาน AUC $0.8939$ \\
    \midrule
    \textbf{Task 2} & $2 \times 2$ Latent Grid + ฟีเจอร์อิสระเชิงมีเงื่อนไข ($|\bar{r}|=0.0036$) + Label Flips $3\%$ & Histogram EM-NB (Inner LL-Select) + Blend + Fixed-Point Affine $F_1$ & LightGBM โอเวอร์ฟิต $+0.0416$ เพราะใช้ขั้นบันไดจำลองเส้นเรียบ 23 มิติ; แตะเพดาน Bayes $\approx 0.758\text{--}0.760$ \\
    \midrule
    \textbf{Task 3} & Zero-Inflated Hurdle + ตัวคูณความชัน 1.0 บน \feat{prev} + Bounded Trapezoid Noise + 266 Whale Spikes & 10-Seed LGBM $P(y>0)$ + 10-Seed Shallow CatBoost \feat{d4} MultiQuantile บน Offset & ต้นไม้ลึกจำ noise; Fitted Offsets ถูก Spike บิดเบือน; MLP ($0.78430$) แพ้ Step Factors; แตะเพดาน $\sim 0.796\text{--}0.797$ \\
    \midrule
    \textbf{Task 4} & \textit{\TFourStatus} & \textit{\TFourStatus} & \textit{สงวนพื้นที่ตามข้อตกลง (หัวข้อย่อย~\ref{subsec:task4})} \\
    \midrule
    \textbf{Task 5} & พื้นหลังปกติอิสระ ($\max |r|=0.0263$) + Temporal Shocks + 1-D Gap \& 3-D Uniform Boxes + NaN & 4-Stage Residual Loop + Exact Support Boxes + NaN-Aware NB ($\text{post}>0.10$) & IForest (42.2\% FP) และ TMAE (82.8\% FP) ถูก 20+ มิติปกติกลบกลุ่มคงที่ 1--3 มิติ; ส่วนเหลือ Null ($z \le 6.27$) \\
    \bottomrule
  \end{tabularx}
\end{table}

\subsection{ข้อคิดในการประยุกต์ใช้ Machine Learning กับอุตสาหกรรมเกม (Insights for the Game Industry)}
\label{subsec:game_industry_insights}

ผลวิเคราะห์จากชุดข้อมูลของ Karena Thailand สะท้อนหลักปฏิบัติสำคัญ 4 ประการสำหรับการนำระบบ AI/ML ไปปฏิบัติงานจริงในอุตสาหกรรมเกมออนไลน์:

\begin{enumerate}
  \item \textbf{แยก ``ระบบคัดกรอง (Triage Queue)'' ออกจาก ``คำสั่งแบนอัตโนมัติ (Auto-Ban)'' เพื่อปกป้องผู้เล่นทักษะสูง:}
  ใน Task 1 ข้อมูลชี้ว่าผู้เล่นใหม่ที่บริสุทธิ์บางกลุ่มมีลักษณะคล้าย ``Smurf'' โดยมีค่าสถิติการยิงสูงกว่าผู้โกงเสียอีก ($\text{KD} = 5.6$, $\text{headshot} = 53\%$ เทียบกับค่าเฉลี่ยผู้โกง $4.46$ และ $45.4\%$) ณ จุดทำงานที่เหมาะกับ Recall ($\text{thr}=0.3769$, $\text{Recall} \approx 97.2\%$) ความแม่นยำรายรายอยู่ที่ราว $54.4\%$ การใช้คะแนนระดับนี้แบนอัตโนมัติจะทำลายผู้เล่นบริสุทธิ์เช่นเดียวกับระบบกฎเดิมที่มี FP $70\%$ แต่เมื่อออกแบบเป็นระบบ \textbf{Risk-Stratified Triage 3 ระดับ} (Tier 1 ปัดตกอัตโนมัติ $40.67\%$, Tier 2 คิวตรวจพฤติกรรมเสริม $41.96\%$, และ Tier 3 ส่งตรงทีม Manual Review $500$ เคส/วัน ที่คะแนน $\ge 0.85$ ซึ่งมี OOF Precision สูงถึง $95.56\%$) จะช่วยระบาย Backlog $45{,}000+$ รายการได้จริงโดยไม่กระทบผู้เล่นสุจริต

  \item \textbf{เศรษฐศาสตร์ของกลุ่ม Whale และการเลือกตัวพยากรณ์ให้ตรงกับฝ่ายธุรกิจ (Finance vs.\ Marketing):}
  ความเหลื่อมล้ำของการใช้จ่ายในเกมสูงมาก ใน Task 3 กลุ่ม Whale ($14.02\%$ ของผู้เล่น) สร้างรายได้ถึง \textbf{87.28\% ของยอดใช้จ่ายรวม} (และใน Task 2 กลุ่มผู้จ่ายสูงสุด $5\%$ สร้างรายได้ $60\%$) แต่การนำแบบจำลองเดียวไปใช้ผิดบริบทจะเกิดความเสียหายร้ายแรง ดังข้อควรระวังในหัวข้อย่อย~\ref{subsubsec:t3_business_impact}: \textbf{Conditional Mixture Median ($\hat{y}_{\text{med}}$)} เหมาะสำหรับฝ่ายการตลาดในการระงับคูปองลด 50\% กับผู้ที่พร้อมจ่ายเต็มราคาอยู่แล้ว (ลดการรั่วไหลกว่า $6$ ล้านบาท/เดือน) \textbf{แต่ห้ามนำไปบวกรวมเพื่อพยากรณ์รายได้องค์กร} เพราะในกลุ่ม Casual ค่าเฉลี่ยของ Median ต่ำกว่ายอดจริงถึง $-93.5\%$ ($4.81$ vs $73.96$ THB) ฝ่ายการเงินจึงต้องใช้ผลรวมของ \textbf{Conditional Expected Revenue} ($\sum \hat{p}_i \widehat{\mathbb{E}}[Y_i \mid Y_i > 0, \mathbf{x}_i]$) ตามหลัก Linearity of Expectation

  \item \textbf{การจัดลำดับความสำคัญด้วยพฤติกรรมร่วม (Multi-Dimensional Persona \& VIP Triage):}
  ตัวชี้วัดเชิงเดี่ยวอย่างชั่วโมงการเล่น (\feat{total\_playtime\_hours}) ไม่สามารถแยก Grinder ($1{,}976.8$ ชม., มัธยฐาน $28{,}611$ บาท) ออกจาก Whale ($2{,}015.1$ ชม., มัธยฐาน $45{,}093$ บาท) ได้ การใช้ความน่าจะเป็นที่ปรับเทียบแม่นยำจาก EM-NB (Log-Loss $0.5975$) ร่วมกับ \feat{p\_whale} จาก Task 3 ช่วยให้ทีมสนับสนุน VIP 20 คนจัดลำดับ Ticket ของผู้เล่นมูลค่าสูงระดับ $255{,}000$ THB ก่อนผู้เล่นสายเล่นนานที่ไม่เติมเงิน ($150$ รายรวมเพียง $12{,}000$ THB) และยกระดับประสิทธิผลแคมเปญการตลาดจาก $2\text{--}5\%$ ได้ตามเป้าหมาย $3\text{--}4$ เท่า

  \item \textbf{ระบบความปลอดภัยต้องให้ ``อนุกรมวิธานของภัยคุกคาม'' ไม่ใช่แค่คะแนนเตือนภัย:}
  แทนที่จะใช้ระบบกฎเดิมที่มี FP สูงถึง $60\%$ หรือโมเดลกล่องดำที่บอกไม่ได้ว่าผู้เล่นผิดปกติเพราะเหตุใด สถาปัตยกรรม 4 ชั้นใน Task 5 สามารถจำแนกบัญชีแจ้งเตือน $2{,}104$ ราย ($8.13\%$) ออกเป็น 8 ประเภทปฏิบัติการเฉพาะทาง เช่น บังคับ 2FA และย้อนธุรกรรมสำหรับ \emph{Account Takeover} ($310$), ระงับคิว Ranked สำหรับ \emph{Account Selling} ($311$), ล็อกสกุลเงินการชำระเงินโดยไม่แบนนักเดินทางสำหรับ \emph{VPN Exploit} ($198$), ส่ง CAPTCHA เชิงพฤติกรรมสำหรับ \emph{Bot Farms} ($207$), และเฝ้าระวังเงียบสำหรับ \emph{Stationary Clusters} ($1{,}070$)
\end{enumerate}

\subsection{แนวทางการพัฒนาต่อยอดในอนาคต (Future Improvements \& Extensions)}
\label{subsec:future_work}

เพื่อยกระดับระบบจากต้นแบบการแข่งขันสู่ระบบวิศวกรรมระดับโปรดักชันของ Karena Thailand ทีมขอเสนอแนวทางการพัฒนาต่อยอด 5 ด้าน ดังนี้:

\begin{enumerate}
  \item \textbf{ระบบเฝ้าระวัง Telemetry Drift และอัตราค่าสูญหาย (Missingness-Rate Monitoring):}
  จากประสบการณ์ใน Task 1 ที่อัตราค่าสูญหายระหว่างชุดฝึกกับชุดทดสอบเปลี่ยนจาก $3.71$ เป็น $2.48$ คอลัมน์/แถว (Adversarial AUC $= 0.90$) ในระบบจริงหลังอัปเดตแพตช์เกม ควรติดตั้งตัวติดตาม Missingness Drift ของแต่ละฟีเจอร์แบบเรียลไทม์ เพื่อแยกระหว่าง ``คุณภาพท่อส่งข้อมูลดีขึ้น'' กับ ``การเสื่อมสภาพของโมเดล'' โดยไม่ต้องปรับจุดตัดตามสถิติชุดทดสอบ

  \item \textbf{การขยายมิติข้อมูลสู่ระดับลำดับเหตุการณ์และเซสชัน (Granular Event \& Sequence Telemetry):}
  ข้อจำกัดของเพดานข้อมูลในปัจจุบันเกิดจากการสรุปพฤติกรรมเป็นตัวเลขรวมรายเดือน (Task 1, Task 3) หรือ 4 ช่วงเวลาหยาบ (Task 5) การนำข้อมูลลำดับเหตุการณ์ระดับแมตช์หรือรายชั่วโมง (เช่น เส้นทางการเล็งรายเฟรม, ลำดับเวลาซื้อไอเทม, หรือกราฟการส่งของขวัญระหว่างบัญชี) มาสร้างฟีเจอร์เชิงเวลาจะช่วยทลายเพดาน Bayes ของฟีเจอร์สถิติแบบเดิมได้

  \item \textbf{สถาปัตยกรรมพยากรณ์รายได้แบบสองหัว (Dual-Head Spend Architecture: Median + Expectation):}
  พัฒนาสายงาน Task 3 จาก \feat{bag\_q10p10} (ซึ่งทำนายควอนไทล์ช่วง $[0.02, 0.50]$ เพื่อ NMAE/การตลาด) ให้มีหัวพยากรณ์คู่ขนานที่ประมาณค่าคาดหมาย $\widehat{\mathbb{E}}[Y \mid Y>0, \mathbf{x}]$ โดยตรง (เช่น อินทิเกรตสเปกตรัมควอนไทล์เต็มช่วง $\alpha \in (0, 1)$ ร่วมกับโมเดลส่วนหาง) สำหรับฝ่ายการเงิน รวมถึงเพิ่มสัญญาณความถี่ธุรกรรมระยะสั้นเพื่อลดความเสี่ยงจากพฤติกรรม Whale Spikes

  \item \textbf{วงจรป้อนกลับแบบ Human-in-the-Loop สำหรับความปลอดภัยและป้องกันการโกง:}
  เชื่อมผลการตรวจสอบบัญชีจริงจากทีมผู้ดูแล (Moderator Decisions ใน T1 Tier 3 และ T5 Stage A) กลับมาเป็นป้ายกำกับแบบ Semi-Supervised / Active Learning เพื่อปรับเทียบจุดตัด $p^*$ และค้นหากลุ่มพฤติกรรมโกงหรือบอทสายพันธุ์ใหม่ที่วิวัฒนาการหนีกฎกล่องเดิมได้อย่างต่อเนื่อง

  \item \textbf{การเสียบต่อส่วน Task 4 เข้าสู่สถาปัตยกรรม LTV สมบูรณ์:}
  เมื่อเกณฑ์การประเมินและมาตรฐานข้อมูลของ Task 4 (Game Title Detection) นิ่งสมบูรณ์แล้ว สามารถนำเนื้อหาและแบบจำลองของ Task 4 มาบรรจุลงในช่องที่สงวนไว้ (\feat{report/sections/sec\_04\_task4\_placeholder.tex}, หัวข้อย่อย~\ref{subsec:task4}) เพื่อเชื่อมต่อระบบคัดแยกภาพหน้าจอเข้าสู่ระบบจัดเส้นทาง Ticket ซัพพอร์ตและตรวจสอบกิจกรรมข้ามเกมทั้ง 5 เกมของ Karena Thailand ได้อย่างสมบูรณ์
\end{enumerate}
